\documentclass{article}
\usepackage{amsmath}
\usepackage{listings}
\usepackage{microtype}
\begin{document}

\section{Projektorganisation und Durchführung}
\label{sec:projektorganisation}
Dieses Kapitel beschreibt die organisatorische Durchführung des Projekts:
die Themenfindung, die Aufteilung der Aufgaben im Team sowie die
Kommunikationsstruktur während der Projektlaufzeit. Im Gegensatz zu den
technischen Kapiteln liegt der Fokus hier auf dem Projektverlauf als
Team, nicht auf einzelnen Implementierungsdetails.

\subsection{Themenfindung}
\label{sec:themenfindung}
Zu Beginn des Projekts wurden im Team mehrere mögliche Themen
vorgeschlagen und in gemeinsamen Meetings diskutiert. Nach Abwägung der vorgeschlagenen 
Alternativen entschied sich das Team für die Entwicklung eines akustischen Ray-Tracers mit 
anschließender digitaler Audioverarbeitung, da dieses Thema sowohl eine substantielle theoretische 
Grundlage (Akustik, Monte-Carlo-Integration,
digitale Signalverarbeitung) als auch einen nicht-trivialen
Implementierungsaufwand (parallele Hochleistungsberechnung in Rust,
Cluster-Verteilung, Audio-Pipeline in Python) vereint.

\subsection{Aufgabenverteilung}
\label{sec:aufgabenverteilung}
Die Aufgabenverteilung im Team erfolgte nicht anhand eines im Voraus
festgelegten Gesamtplans, sondern iterativ und abhängig vom
Projektfortschritt. Da einzelne Arbeitspakete inhaltlich aufeinander
aufbauten (z.\,B. konnte die Definition des JSON-Datenschemas erst nach
Festlegung der von der Physics Engine gelieferten Werte erfolgen, und die
DSP-Pipeline setzte wiederum ein funktionierendes Datenschema voraus),
ergaben sich viele konkrete Arbeitspakete erst im Projektverlauf. Zwei
Teammitglieder erklärten sich früh bereit, sich mit der Rust-Entwicklung
der Physics Engine auseinanderzusetzen; die übrigen Aufgaben -- Audio-DSP,
Cluster-Infrastruktur, Geometrie-Import, Visualisierung sowie die
Dokumentation -- wurden im weiteren Verlauf sukzessive definiert und
zugeteilt, sobald die jeweiligen Vorgänger-Arbeiten einen ausreichenden
Stand erreicht hatten.

Beispiele für auf diese Weise entstandene Arbeitspakete sind unter
anderem:
\begin{itemize}
 \item Recherche realistischer Absorptions- und Streukoeffizienten für
 Materialien sowie Definition des zugehörigen JSON-Schemas,
 \item Implementierung der Serialisierung der Simulationsergebnisse
 (\texttt{serde}/\texttt{serde\_json}) in ein einheitliches
 \texttt{ir\_output.json}-Format,
 \item Anbindung von Python und Rust über einen Subprocess-Aufruf
 inklusive JSON-Ingestion,
 \item Erweiterung der Faltungslogik um frequenzabhängige,
 materialspezifische Dämpfung,
 \item Umstellung der Geometrieverarbeitung von 2D auf 3D
 (Möller--Trumbore-Algorithmus) sowie Anpassung der Visualisierung
 auf einen echten 3D-Plot,
 \item Aufbau und Absicherung der Cluster-Infrastruktur zur verteilten
 Ausführung.
\end{itemize}

\subsection{Kommunikation}
\label{sec:kommunikation}
Die Abstimmung im Team erfolgte über regelmäßige
Meetings, in denen der Gesamtfortschritt besprochen, offene Fragen
geklärt und anstehende Arbeitspakete verteilt wurden, sowie ein
fortlaufender Austausch über einen Discord-Kanal zwischen den Meetings.
Fertiggestellte Arbeitspakete wurden im Team gegenseitig gesichtet, bevor
mit darauf aufbauenden Aufgaben begonnen wurde.

\subsection{Herausforderungen im Projektverlauf}
\label{sec:herausforderungen-orga}
Eine zentrale organisatorische Herausforderung während der
Projektlaufzeit war der Ausstieg des ursprünglichen Teammitglieds, der unter
anderem für den Aufbau der LaTeX-Dokumentation zuständig war. Dessen Aufgaben mussten 
kurzfristig auf die verbleibenden Teammitglieder umverteilt werden, was zu einer temporären 
Verzögerung im Projektfortschritt führte. Als Team-Prinzip
wurde daher etabliert, dass alle Mitglieder bei Bedarf gegenseitig
einspringen und sich nicht ausschließlich auf das eigene zugewiesene
Arbeitspaket beschränken. Die Gesamtdokumentation wird
als gemeinsames Werk des Teams verstanden und nicht nach Autoren
einzelner Abschnitte kenntlich gemacht, sodass alle Mitglieder gemeinsam
für Konsistenz, Vollständigkeit und Korrektheit des Gesamtdokuments
verantwortlich sind.

\end{document}
